% Generated draft: requires mathematical and source review.
\documentclass[11pt]{article}
\usepackage[margin=25mm]{geometry}
\usepackage{fontspec}
\setmainfont{Latin Modern Roman}
\usepackage{amsmath,amssymb,mathtools,unicode-math}
\setmathfont{Latin Modern Math}
\usepackage{enumitem,needspace,xurl,hyperref}
\setlength{\parindent}{0pt}
\setlength{\parskip}{5pt}
\setlength{\emergencystretch}{3em}
\newcommand{\N}{\mathbb{N}}
\newcommand{\Z}{\mathbb{Z}}
\newcommand{\Q}{\mathbb{Q}}
\newcommand{\R}{\mathbb{R}}
\newcommand{\C}{\mathbb{C}}
\newcommand{\F}{\mathbb{F}}
\newcommand{\A}{\mathbb{A}}
\newcommand{\PP}{\mathbb{P}}
\newcommand{\E}{\mathbb{E}}
\newcommand{\eps}{\varepsilon}
\newcommand{\dd}{\,\mathrm{d}}
\begin{document}
\section*{UnsolvedMath ID: 30000514}
\textbf{Source problem code:} catalogue entry
\subsection{Definitions and mathematical statement}
Let \(n,m,d\) be integers satisfying
\[
 1\leq m\leq n-3,\qquad d\geq2n+2m-1,
\]
and let \(\Gamma\subset\mathbb P^{n-1}\) be a set of \(d\) distinct points in the source's symmetric-position sense. The ground field and that symmetry condition are not specified. For points over a field \(k\), a standard Hilbert-function convention is
\[
 h_\Gamma(t)=\dim_k(k[x_0,\ldots,x_{n-1}]/I_\Gamma)_t,
\]
where \(I_\Gamma\) is the homogeneous vanishing ideal; equivalently, it is the number of independent evaluations imposed in degree \(t\), with the appropriate field and homogeneous-fiber interpretation. If geometric points are used, the field of scalar extension must be fixed consistently.

The assertion as printed is
\[
 h_\Gamma(2)\geq2n+m-2
 \quad\Longrightarrow\quad
 \exists C\subset\mathbb P^{n-1}\text{ a curve},\quad
 \Gamma\subseteq C,\quad\deg C\leq n+m-2,
\]
for every \(\Gamma,n,m,d\) satisfying the stated hypotheses. The Hilbert-function condition is a lower bound in both the original and cleaned records. It is retained exactly as \(\geq\), not silently changed to \(\leq\) or equality. Confirmation of its intended sign is particularly important for identifying which quantitative Castelnuovo-type statement the references address.

Symmetric position can in an arithmetic setup refer to the full symmetric group on the geometric points being induced by a Galois action, or in a family to a specified monodromy group. A weaker uniform-position statement about Hilbert functions of subsets is not automatically the same condition. No Galois field, family or precise symmetry predicate is supplied, so none of these is selected as the definition. There is also no explicit linear-general-position or nondegeneracy assumption to add.

The source does not say whether the curve must be integral, reduced, irreducible or pure of dimension one, or defined over the original field rather than an extension. Degree is conventionally its projective degree, the leading coefficient of a curve's Hilbert polynomial. The admissible curve convention matters, especially if arbitrary zero-dimensional components were allowed, and needs confirmation.

The literature triage checked 2026-08-21 reports proofs for the cases \(m=3,4\) and says the stated range is not fully settled. The full range \(1\leq m\leq n-3\), including \(m=1,2\), is retained; the description of \(3,4\) as first cases does not remove those values. The cited theorem hypotheses and compatibility with the printed inequality are not reproduced. The partially-solved label and verification report remain supplied metadata.
\subsection{Short English statement}
Under the stated point-count and symmetric-position assumptions, does the printed Hilbert-function lower bound force containment in a curve of degree at most \(n+m-2\)? The inequality sign and full parameter range are preserved. The field, symmetry condition and allowed type of curve are missing and must be checked before matching the reported special cases.
\subsection{Sources}
\begin{itemize}
\item Local UnsolvedMath registry, permanent ID 30000514, source code catalogue entry, ``Low-Degree Curves Through Points in Symmetric Position,'' including both statements with \(h_\Gamma(2)\geq2n+m-2\) and verification metadata.
\item \emph{Classical Algebraic Geometry (2007)}, Oberwolfach Reports, pp.\ 1615--1662, DOI 10.4171/owr/2006/27: \url{https://doi.org/10.4171/owr/2006/27}. The title uses 2007, while the DOI and proposed-year metadata use 2006.
\item Supplied reference ``Petrakiev, Castelnuovo theory via Gr\"obner bases'': \url{https://arxiv.org/abs/math/0604517}.
\item Supplied reference ``higher Castelnuovo theory'': \url{https://www.numdam.org/item/AST_1993__218__187_0/}.
\item Literature triage checked 2026-08-21; record updated 2026-09-16. No external source has been visited or independently checked here.
\end{itemize}
\end{document}
